\subsection{Produktkrav}
\label{sec:produktkrav}
\subsubsection{Kravgrunnlag og prioritering}
Kravspesifikasjonen skiller mellom rammene i oppgavebeskrivelsen (R1), som er obligatoriske, og egne mål for ytelse og konstruksjon (R2--R8). Akseptgrensene for R3--R6 ble satt etter den første kjøretesten.

\subsubsection{Kravspesifikasjon}
\begin{description}
    \item[R1: Oppgitte rammer.] Ytre mål skal ikke overstige $140 \times 85 \times 70$ mm. Det utdelte stålsvinghjulet (Ø60 × 6 mm) og den utdelte akselen (Ø6 × 65 mm) skal brukes, og bilen trekkes opp med 20 cm tråd. Verifiseres ved måling og inspeksjon mot oppgavegrunnlaget.
    \item[R2: Funksjon.] Bilen skal kjøre selvstendig etter opptrekk. Verifiseres ved funksjonstest.
    \item[R3: Hovedmål.] Størst mulig gjennomsnittlig kjøredistanse ved standardisert opptrekk, og minst \SI{5}{m}. Vurderes med minst 10 kjøringer per sammenlignet konfigurasjon.
    \item[R4: Repeterbarhet.] Standardavviket for kjøredistansen er under \SI{15}{\percent} av snittet.
    \item[R5: Retningsstabilitet.] Sideavviket er i snitt under \SI{10}{\percent} av lengden, målt etter prosedyren i seksjon~\ref{sec:metode-kjoretest}.
    \item[R6: Pålitelighet.] Minst 9 av 10 opptrekk gir en vellykket kjøring, det vil si at bilen kjører av seg selv etter nedsettingen og stopper uten å treffe noe.
    \item[R7: Justerbarhet.] Drivverket bør kunne demonteres og justeres, og utveksling og hjul bør kunne endres uten fullstendig redesign. Verifiseres ved inspeksjon og praktisk utprøving.
    \item[R8: Produksjon.] Bilen skal kunne fremstilles med tilgjengelige produksjonsmetoder. Verifiseres mot produksjonsunderlag og ferdig prototype.
\end{description}

